% !TEX root = ../main.tex
\chapter{ケーススタディ：口座操作の仕様証明}
\label{chap:account-spec-proof}

\begin{chapterabstract}
本章では、入金、出金、送金という小さな口座操作を題材に、自然言語の仕様を Lean 4 で検査できる形へ落とし込む。中心になるのは、残高が負にならないこと、出金成功時に指定額だけ残高が減ること、出金失敗時に状態が変わらないこと、送金成功時に二つの口座の総残高が保存されることである。現実の金融システム全体を検証するのではなく、仕様の核を小さく切り出し、型、関数、述語、定理として表す方法を学ぶ。
\end{chapterabstract}

\begin{learninggoals}
\item 口座操作の自然言語仕様を、状態、操作、事前条件、事後条件、不変条件へ分解できる。
\item \texttt{Account}、\texttt{Balance}、\texttt{Amount} を Lean 4 の小さなモデルとして定義できる。
\item \texttt{deposit}、\texttt{withdraw?}、\texttt{transfer?} の仕様を等式と場合分けで証明できる。
\item Lean で証明した性質と、実システム側で別途確認すべき性質を区別できる。
\end{learninggoals}

\section{この章を読む理由}

口座操作は、仕様証明の入門例として扱いやすい。入金、出金、送金という操作は直感的でありながら、状態更新、不変条件、事前条件、成功ケース、失敗ケースをすべて含んでいるからである。

たとえば、次のような仕様はよく見かける。

\begin{quote}
出金額が残高以下なら出金は成功し、残高は出金額だけ減る。出金額が残高を超えるなら出金は失敗し、残高は変わらない。送金は、送金元から出金し、送金先へ入金する。送金に成功した場合、二つの口座の総残高は変わらない。
\end{quote}

この文章は人間には自然に読める。しかし、機械に検査させるには、いくつかの部品へ分ける必要がある。どの型が状態なのか。どの関数が操作なのか。どの条件が成功条件なのか。どの等式を証明すれば、仕様の重要部分を押さえたと言えるのか。

\begin{intuitionbox}
口座操作の仕様証明とは、「たまたまいくつかのテストで合う」ことを確認する作業ではない。すべての口座、すべての金額について、モデル化した範囲の仕様が必ず成り立つことを Lean に確認させる作業である。
\end{intuitionbox}

\FigurePlaceholder{fig-ch19-account-flow.png}{0.92}{入金・出金・送金を状態遷移として見る}{fig:ch19-account-flow}

図\ref{fig:ch19-account-flow}では、口座を状態、入金・出金・送金を状態遷移として見る。出金には成功と失敗の分岐があり、送金は出金と入金を組み合わせた操作として読む。

\section{自然言語仕様を分解する}

最初に、仕様を Lean に移しやすい粒度へ分解する。本章では、次のように整理する。

\begin{longtable}{p{0.22\linewidth}p{0.36\linewidth}p{0.32\linewidth}}
\toprule
分類 & 本章での意味 & Lean での表現 \\
\midrule
状態 & 口座IDと残高を持つ口座 & \texttt{Account} \\
値 & 残高と操作金額 & \texttt{Balance}, \texttt{Amount} \\
操作 & 入金、出金、送金 & \texttt{deposit}, \texttt{withdraw?}, \texttt{transfer?} \\
事前条件 & 出金額が残高以下であること & \texttt{CanWithdraw a amount} \\
事後条件 & 成功時・失敗時に結果がどうなるか & 等式としての \texttt{theorem} \\
不変条件 & 残高が負にならないこと、総残高が保存されること & 述語と保存定理 \\
\bottomrule
\end{longtable}

\begin{softwarebox}
ここで作る Lean モデルは、本番コードの完全な置き換えではない。設計レビューで「この操作はどの性質を守るべきか」を明確にするための小さな検証モデルである。
\end{softwarebox}

\subsection{証明する性質を選ぶ}

本章で証明する性質は、次の四つに絞る。

\begin{enumerate}
\item 入金後の残高は、元の残高に入金額を足した値である。
\item 出金可能な場合、出金結果は成功し、残高は指定額だけ減る。
\item 出金不可能な場合、出金結果は失敗する。
\item 送金に成功した場合、送金元と送金先の総残高は保存される。
\end{enumerate}

これだけでも、仕様の重要な核はかなり明確になる。特に四つ目の「総残高が保存される」は、送金という複合操作が単なる二つの更新ではなく、全体として保存すべき量を持つことを表している。

\begin{warningbox}
本章では、通貨単位、小数、丸め、手数料、監査ログ、同時実行、永続化、認可、外部API呼び出しは扱わない。これらは実務上は重要だが、最初の Lean モデルにすべて入れると、証明したい核が見えにくくなる。
\end{warningbox}

\section{ドメインモデルを作る}

本章では、残高と金額を自然数 \texttt{Nat} で表す。自然数は負にならない数である。そのため、「残高は負にならない」という不変条件の一部を型の選択によって表している。

\begin{definitionbox}
\textbf{不変条件}とは、操作の前後で常に保たれるべき条件である。本章では、残高が自然数であることにより、残高非負条件を型のレベルで表す。さらに、送金では二つの口座の総残高が保存されることを定理として表す。
\end{definitionbox}

\begin{leanbox}
\texttt{abbrev} は既存の型に別名を与える。ここでは \texttt{Balance} と \texttt{Amount} をどちらも \texttt{Nat} の別名としておく。これにより、Lean 上では同じ自然数でも、本文では「残高」と「操作金額」を区別して読める。
\end{leanbox}

\begin{listing}[htbp]
\begin{minted}{lean4}
namespace Chapter19

abbrev AccountId := Nat
abbrev Balance := Nat
abbrev Amount := Nat

structure Account where
  id : AccountId
  balance : Balance
deriving Repr, DecidableEq

def NonNegative (a : Account) : Prop :=
  0 <= a.balance

theorem account_nonnegative (a : Account) : NonNegative a := by
  exact Nat.zero_le a.balance
\end{minted}
\caption{口座操作の最小ドメインモデル}
\label{lst:ch19-domain-model}
\end{listing}

このモデルでは、\texttt{Account} は口座IDと残高だけを持つ。現実には、凍結状態、通貨、所有者、更新時刻なども必要になるかもしれない。しかし本章では、残高操作の仕様を見たいので、余計なフィールドは入れない。

\subsection{型に入れた仕様と、定理で証明する仕様}

\texttt{Balance := Nat} とした時点で、残高が負にならないことはほとんど型によって保証される。一方で、「出金成功時は指定額だけ減る」や「送金では総残高が保存される」は、型だけでは表せない。これらは関数の振る舞いに関する主張なので、定理として証明する。

\begin{softwarebox}
型で表せる制約は型に入れる。操作の振る舞いに関する制約は定理で表す。この切り分けを意識すると、Lean モデルは読みやすくなる。
\end{softwarebox}

\section{操作を関数として定義する}

次に、入金、出金、送金を関数として定義する。

入金は常に成功する操作として扱う。出金は、金額が残高以下なら成功し、そうでなければ失敗する。失敗を表すために、ここでは \texttt{Option Account} を使う。\texttt{some a} は成功して新しい口座 \texttt{a} が得られたことを表し、\texttt{none} は失敗を表す。

\begin{listing}[htbp]
\begin{minted}{lean4}
def deposit (a : Account) (amount : Amount) : Account :=
  { a with balance := a.balance + amount }

def withdraw? (a : Account) (amount : Amount) : Option Account :=
  if amount <= a.balance then
    some { a with balance := a.balance - amount }
  else
    none

def totalBalance (a b : Account) : Balance :=
  a.balance + b.balance

def transfer? (src to : Account) (amount : Amount) :
    Option (Account × Account) :=
  match withdraw? src amount with
  | none => none
  | some src' => some (src', deposit to amount)
\end{minted}
\caption{入金・出金・送金の定義}
\label{lst:ch19-operations}
\end{listing}

\texttt{transfer?} は、まず送金元から出金しようとする。出金に失敗した場合、送金全体も失敗する。出金に成功した場合、送金先へ同じ金額を入金し、更新後の二つの口座を返す。

\begin{intuitionbox}
送金は「出金してから入金する」処理の合成である。ただし、途中の出金が失敗する可能性があるため、普通の関数合成ではなく、失敗を含む合成として読む。
\end{intuitionbox}

\FigurePlaceholder{fig-ch19-withdraw-transfer.png}{0.92}{出金の成功・失敗と送金の合成}{fig:ch19-withdraw-transfer}

\section{入金の事後条件を証明する}

まず、最も簡単な仕様から証明する。入金後の残高は、元の残高に入金額を足した値である。

\begin{listing}[htbp]
\begin{minted}{lean4}
theorem deposit_balance (a : Account) (amount : Amount) :
    (deposit a amount).balance = a.balance + amount := by
  rfl

theorem deposit_preserves_nonnegative
    (a : Account) (amount : Amount) :
    NonNegative (deposit a amount) := by
  exact Nat.zero_le (deposit a amount).balance
\end{minted}
\caption{入金後の残高仕様}
\label{lst:ch19-deposit-proof}
\end{listing}

\texttt{deposit\_balance} は \texttt{rfl} で証明できる。\texttt{rfl} は、定義を展開すれば左右が同じであるときに使える。

\begin{definitionbox}
\texttt{rfl} は「反射律」による証明である。Lean が定義を展開した結果、左辺と右辺が同じ形になる場合、\texttt{rfl} で証明できる。
\end{definitionbox}

ここで大事なのは、証明が短いことではない。仕様が \texttt{deposit\_balance} という名前で残ることである。この名前は設計文書やレビューコメントから参照できる。

\section{出金の事前条件と事後条件を証明する}

出金には成功と失敗の分岐がある。そのため、まず成功条件を名前付きの述語として定義する。

\begin{listing}[htbp]
\begin{minted}{lean4}
def CanWithdraw (a : Account) (amount : Amount) : Prop :=
  amount <= a.balance

theorem withdraw_success (a : Account) (amount : Amount)
    (h : CanWithdraw a amount) :
    withdraw? a amount =
      some { a with balance := a.balance - amount } := by
  have h' : amount <= a.balance := h
  simp [withdraw?, h']

theorem withdraw_failure (a : Account) (amount : Amount)
    (h : ¬ CanWithdraw a amount) :
    withdraw? a amount = none := by
  have h' : ¬ amount <= a.balance := by
    intro hp
    exact h hp
  simp [withdraw?, h']
\end{minted}
\caption{出金可能条件と出金仕様}
\label{lst:ch19-withdraw-proof}
\end{listing}

\texttt{withdraw\_success} は、「出金額が残高以下」という証拠 \texttt{h} があるなら、\texttt{withdraw?} は \texttt{some} を返し、その残高は \texttt{a.balance - amount} になる、という主張である。

\texttt{withdraw\_failure} はその逆である。出金可能条件が成り立たないなら、結果は \texttt{none} になる。

\begin{leanbox}
ここで使っている \texttt{simp} は、\texttt{if} の条件が真か偽か分かっている場合に、該当する分岐へ式を整理する。成功ケースでは \texttt{h'} により \texttt{if} の then 側へ進み、失敗ケースでは else 側へ進む。
\end{leanbox}

\subsection{失敗時に状態を変えない設計}

本章の \texttt{withdraw?} は、失敗した場合に新しい口座を返さない。これは「失敗時には状態を更新しない」という設計を \texttt{Option} で表している。

実装言語では、例外、戻り値、トランザクション、イベントログなどで同じ意図を表すことがある。Lean モデルでは、まず \texttt{none} によって「更新後状態が存在しない」と表すことで、仕様を小さく保つ。

\begin{warningbox}
\texttt{withdraw?} が \texttt{none} を返すことは、「本番DBが絶対に更新されない」ことを直接証明しているわけではない。本章で証明しているのは、Lean 上のモデル化された関数について、失敗時に新しい口座状態を返さないという性質である。
\end{warningbox}

\section{送金の仕様を証明する}

送金は、出金と入金を組み合わせた操作である。そのため、出金可能なら送金は成功し、出金不可能なら送金は失敗する。

\begin{listing}[htbp]
\begin{minted}{lean4}
theorem transfer_success (src to : Account) (amount : Amount)
    (h : CanWithdraw src amount) :
    transfer? src to amount =
      some ({ src with balance := src.balance - amount },
            { to with balance := to.balance + amount }) := by
  have h' : amount <= src.balance := h
  simp [transfer?, withdraw?, deposit, h']

theorem transfer_failure (src to : Account) (amount : Amount)
    (h : ¬ CanWithdraw src amount) :
    transfer? src to amount = none := by
  have h' : ¬ amount <= src.balance := by
    intro hp
    exact h hp
  simp [transfer?, withdraw?, h']
\end{minted}
\caption{送金の成功・失敗仕様}
\label{lst:ch19-transfer-cases}
\end{listing}

ここまでで、送金の成功条件と失敗条件は明確になった。次に、成功した送金が総残高を保存することを証明する。

\subsection{総残高保存の証明方針}

証明したい等式は、次の形である。

\[
  (\mathrm{from.balance} - \mathrm{amount})
  + (\mathrm{to.balance} + \mathrm{amount})
  =
  \mathrm{from.balance} + \mathrm{to.balance}
\]

この等式は、「送金元から減った金額」と「送金先へ増えた金額」が打ち消し合うことを表している。ただし、自然数の引き算は、出金額が残高以下であるという条件があって初めて期待通りに振る舞う。その条件が \texttt{CanWithdraw from amount} である。

\begin{listing}[htbp]
\begin{minted}{lean4}
theorem transfer_success_preserves_total
    (src to : Account) (amount : Amount)
    (h : CanWithdraw src amount) :
    totalBalance { src with balance := src.balance - amount }
      { to with balance := to.balance + amount } =
    totalBalance src to := by
  have h' : amount <= src.balance := h
  dsimp [totalBalance]
  calc
    (src.balance - amount) + (to.balance + amount)
        = (src.balance - amount) + (amount + to.balance) := by
            rw [Nat.add_comm to.balance amount]
    _ = ((src.balance - amount) + amount) + to.balance := by
            rw [← Nat.add_assoc]
    _ = src.balance + to.balance := by
            rw [Nat.sub_add_cancel h']
\end{minted}
\caption{送金成功時の総残高保存}
\label{lst:ch19-transfer-total-proof}
\end{listing}

この証明では、\texttt{calc} を使って式変形の道筋を明示している。まず送金先側の \texttt{to.balance + amount} の順序を入れ替える。次に括弧を付け替える。最後に、\texttt{amount <= from.balance} という条件を使って、\texttt{from.balance - amount + amount = from.balance} を適用する。

\begin{definitionbox}
\texttt{calc} は、等式変形を段階的に書くための構文である。レビューする人は、証明を一つの黒箱としてではなく、どの式変形を行ったかのログとして読める。
\end{definitionbox}

\FigurePlaceholder{fig-ch19-total-preservation.png}{0.92}{送金成功時に総残高が保存される理由}{fig:ch19-total-preservation}

\subsection{結果から総残高保存を読む}

実務上は、\texttt{transfer?} の結果として得られた二つの口座について、総残高が保存されていることを直接使いたい場合がある。その場合は、「\texttt{transfer?} が \texttt{some (from', to')} を返したなら、総残高が保存される」という形の定理にしておくと使いやすい。

\begin{listing}[htbp]
\begin{minted}{lean4}
theorem transfer_result_preserves_total
    (src to src' to' : Account) (amount : Amount)
    (hTransfer : transfer? src to amount = some (src', to')) :
    totalBalance src' to' = totalBalance src to := by
  unfold transfer? at hTransfer
  by_cases h : amount <= src.balance
  · simp [withdraw?, h, deposit] at hTransfer
    rw [← hTransfer.left, ← hTransfer.right]
    exact transfer_success_preserves_total src to amount h
  · simp [withdraw?, h] at hTransfer
\end{minted}
\caption{送金結果から総残高保存を取り出す}
\label{lst:ch19-transfer-result-total}
\end{listing}

この定理では、\texttt{by\_cases} によって、出金可能な場合と不可能な場合に分けている。出金可能な場合は、成功結果の形から総残高保存を使う。出金不可能な場合は、\texttt{transfer?} が \texttt{some} を返したという仮定と矛盾するため、そのケースは消える。

\begin{leanbox}
\texttt{cases hTransfer} は、\texttt{some (...)} 同士が等しいなら中身も等しい、という情報を使って、得られた結果の形を具体化する。これは、API レスポンスやバリデーション結果の成功ケースを分解するときにも使う。
\end{leanbox}

\section{この章の証明を実行例で読む}

Lean の証明は、実行例と組み合わせると読みやすくなる。次の例では、残高100の口座から30を出金し、残高40の口座へ30を送金する。

\begin{listing}[htbp]
\begin{minted}{lean4}
def alice : Account := { id := 1, balance := 100 }
def bob : Account := { id := 2, balance := 40 }

#eval deposit alice 20
#eval withdraw? alice 30
#eval withdraw? alice 130
#eval transfer? alice bob 30
\end{minted}
\caption{実行例で仕様を読む}
\label{lst:ch19-eval-examples}
\end{listing}

\texttt{\#eval} は具体例を確認するためには便利である。しかし、\texttt{\#eval} はあくまで特定の入力に対する実行結果である。\texttt{theorem} は、任意の入力について成り立つ性質を確認する。

\begin{softwarebox}
テストと証明は対立しない。\texttt{\#eval} や通常のテストは、具体例で仕様の読み方を確認するために役立つ。Lean の定理は、モデル化した範囲のすべての入力について性質を保証するために使う。
\end{softwarebox}

\section{仕様レビューでの使い方}

本章の Lean コードは、そのまま本番コードへ貼り付けるためのものではない。仕様レビューで重要なのは、次の対応を確認することである。

\begin{longtable}{p{0.30\linewidth}p{0.30\linewidth}p{0.30\linewidth}}
\toprule
レビュー観点 & Lean 側の対応 & 確認したいこと \\
\midrule
残高型 & \texttt{Balance := Nat} & 負の残高をモデル上排除しているか \\
出金成功条件 & \texttt{CanWithdraw} & 成功条件が仕様書と一致しているか \\
出金成功時の更新 & \texttt{withdraw\_success} & 指定額だけ減ることを示しているか \\
出金失敗時の扱い & \texttt{withdraw\_failure} & 失敗を明示的に表しているか \\
送金成功時の保存量 & \begin{tabular}[t]{@{}l@{}}\texttt{transfer\_success}\\\texttt{\_preserves\_total}\end{tabular} & 総残高保存を示しているか \\
証明済み範囲 & 各 \texttt{theorem} の前提 & 実システムとの差分を説明できるか \\
\bottomrule
\end{longtable}

この表は、設計文書にそのまま近い形で載せられる。Lean の定理名を仕様項目に対応させると、どの性質が証明済みで、どの性質が未証明かを追跡しやすくなる。

\section{テストケースとの棲み分け}

Lean で総残高保存を証明したからといって、送金機能のテストが不要になるわけではない。Lean モデルは、本番システムの一部の性質を抽象化して表したものである。

\subsection{Lean で証明するとよいこと}

\begin{itemize}
\item 金額更新のような小さく純粋なロジック。
\item 出金可能条件と成功・失敗結果の対応。
\item 操作後に保存される数量。
\item リファクタリングしても変わってはいけない等式。
\end{itemize}

\subsection{通常のテストや監視で確認すべきこと}

\begin{itemize}
\item DB トランザクションが正しく張られていること。
\item 並行実行時に二重出金が起きないこと。
\item 外部決済APIとの通信失敗時のリトライ挙動。
\item 監査ログ、認可、通知、UI 表示が仕様通りであること。
\item 本番実装と Lean モデルが乖離していないこと。
\end{itemize}

\begin{warningbox}
「Lean で証明済み」と言うときは、必ず「どのモデルについて、どの前提で、どの定理を証明したのか」を併記する。本章で証明したのは、自然数残高を持つ単純な口座モデルにおける、入金・出金・送金の一部の性質である。
\end{warningbox}

\section{本章のまとめ}

口座操作は、状態、操作、事前条件、事後条件、不変条件を学ぶための小さく有用な題材である。

\texttt{Balance} と \texttt{Amount} を \texttt{Nat} で表すことで、残高非負条件の一部を型に入れた。

\texttt{deposit}、\texttt{withdraw?}、\texttt{transfer?} を関数として定義し、成功・失敗の仕様を等式として証明した。

送金成功時の総残高保存は、\texttt{calc} による式変形としてレビュー可能な形で証明できる。

Lean の証明はテストを置き換えるものではなく、テストでは網羅しにくい仕様の核を補強するものである。
